\begin{table}[H]
\centering
\caption{Classifier-risk robustness: high-confidence composition tests}
\label{tab:appendix_highconf_comp}
\small
\resizebox{\textwidth}{!}{%
\begin{tabular}{>{\raggedright\arraybackslash}p{0.34\textwidth}ccccc}
\hline
Variable & \multicolumn{5}{c}{Dependent variable: log(1 + AI patents)} \\
 & t-2 & t-1 & t & t+1 & t+2 \\
 & (1) & (2) & (3) & (4) & (5) \\
\hline
\multicolumn{6}{l}{\textit{Panel A. Actionable disclosure}} \\
Actionable disclosure (dummy) & 0.047 & 0.049* & 0.030 & 0.035 & 0.072 \\
 & (0.029) & (0.029) & (0.029) & (0.039) & (0.047) \\
\multicolumn{6}{l}{\textit{Panel B. Speculative-only disclosure}} \\
Speculative-only disclosure (dummy) & -0.014 & -0.010 & -0.006 & -0.076 & -0.129 \\
 & (0.049) & (0.046) & (0.044) & (0.080) & (0.082) \\
Controls & Y & Y & Y & Y & Y \\
Firm FE & Y & Y & Y & Y & Y \\
Year FE & Y & Y & Y & Y & Y \\
Adj. $R^2$ & -0.953 & -0.954 & -0.928 & -0.922 & -0.916 \\
Observations & 2,810 & 2,810 & 2,810 & 2,810 & 2,810 \\
\hline
\end{tabular}%
}
\caption*{\footnotesize \textit{Note.} This table reports composition-based robustness checks using only sentences the classifier labels with high confidence. The dependent variable is $\log(1+\text{AI patents})$, and the five columns report the relation at event times $t-2$, $t-1$, $t$, $t+1$, and $t+2$ relative to the disclosure year, so the columns trace how the disclosure indicator lines up with past, current, and future AI patenting. The purpose is to check whether the main patterns hold when the disclosure measures are restricted to the most confidently classified sentences. *, **, and *** denote significance at the 10\%, 5\%, and 1\% levels, respectively.}
\end{table}
